\begin{annexitem}{annex:p3.rn.type-realization}{thm:p03-rn-type-realization}{Compactness with named coordinates and a countable model}
\begin{proof}
Add distinct fresh symbols \(d_0,\ldots,d_{n-1}\), whose
interpretations need not be distinct.
Consistency of \(p\) makes
\[
 T\cup\{\varphi(d_0,\ldots,d_{n-1}):\varphi\in p\}
\]
finitely satisfiable. Compactness gives a model of the full
set. The countable elementary-submodel construction in
the expanded language retains the interpretations of these
constants; the result is infinite because it is a \(T\)-model.
Discarding the new symbols gives the required realization.
Equivalently use the supplied simultaneous construction
with this satisfiable named theory and no omission
requirements (\cref{thm:p03-vaught-simultaneous-construction}).
For \(n=0\) no names are added: the empty tuple realizes
the complete theory's unique zero-type.

For functions, use the total single-valued graph presentation
proved in \cref{thm:p03-vaught-functional-bridge}.
Its relational language is still countable; the model
correspondence preserves domains, formula truth and
the realized tuple type. Apply the preceding construction
there and translate back. The result has the same
countably infinite cardinality.
\end{proof}
\end{annexitem}

\begin{annexitem}{annex:p3.rn.isolated-extension}{thm:p03-rn-isolated-extension}{Existential projection of a consistent isolator}
\begin{proof}
Let \(q\) be the first-\(n\)-slot restriction of \(p\).
Because \(\theta\in p\), the formula
\(\exists y\,\theta(\bar x,y)\) belongs to \(q\).
As \(\bar a\) realizes \(q\), some \(b\) in the same model
satisfies \(\theta(\bar a,b)\).
For each \(\varphi\in p\), isolation gives
\[
 T\models\forall\bar x\forall y
               (\theta(\bar x,y)\to\varphi(\bar x,y)).
\]
Thus \((\bar a,b)\) realizes every formula of \(p\),
and completeness identifies its type with \(p\).
For \(n=0\), the existential projection is a sentence in \(T\),
so this case is included.

Name a finite parameter tuple \(\bar a\), allowing repetitions,
and let \(r(\bar y)\) be a complete type consistent with
the full theory of \((\mathcal M,\bar a)\).
Compactness gives a model of that named theory with a
realization of \(r\).
The joint type \(p(\bar x,\bar y)\) of its named tuple
and realization is a complete empty-parameter type of \(T\).
Its first-coordinate restriction is
\(\operatorname{tp}^{\mathcal M}(\bar a)\).
By hypothesis \(p\) is isolated.
Its existential projection to \(\bar x\) lies in that
restriction, so the preceding argument, now with a finite
tuple \(\bar y\), realizes \(p\) in \(\mathcal M\).
That tuple realizes \(r\) over the original parameters.
Every finite parameter set can be enumerated by \(\bar a\),
which proves exactly finite-parameter saturation.
\end{proof}
\end{annexitem}

\begin{annexitem}{annex:p3.rn.characterization}{thm:p03-rn-characterization}{The omission obstruction and complete finite-parameter back-and-forth}
\begin{proof}
First use the relational/constant setting of Vaught's
construction. Its supplied formula-algebra identification
(\cref{thm:p03-vaught-algebra-types}) identifies complete
types with ultrafilters and isolated types with principal
ultrafilters.
The proved formula-algebra categoricity theorem
(\cref{thm:p03-vaught-finite-formulas-categoricity})
supplies the equivalence of clauses~1--4.
Here are its exact model-theoretic interfaces and the
additional saturation implications; this does not assert
a different original author's proof.

For a nonisolated complete \(n\)-type \(p\), the supplied
simultaneous construction with its one omission requirement
gives a countably infinite model omitting \(p\).
The model of \cref{thm:p03-rn-type-realization} realizes it.
Isomorphisms preserve all formulas on all finite tuples,
so these models cannot be isomorphic.
Clause~1 therefore forces isolation of each complete
finite-tuple type. The unique zero-type is isolated by
\(\top\).
Conversely, isolation of all types makes both countable
models atomic, and the fully proved atomic back-and-forth
(\cref{thm:p03-vaught-atomic-uniqueness}) supplies their
isomorphism.
The complete Boolean-algebra proof
(\cref{thm:p03-vaught-nonprincipal-algebra}) shows that
an infinite formula quotient has a nonprincipal
ultrafilter. Thus isolation of every type forces finiteness
of the formula quotient, not just density of isolated
points. A finite Boolean algebra has finitely many atoms,
every ultrafilter selects exactly one, and all its elements
are unions of atoms. This gives finitely many types and
their isolators. These are the previously supplied
full algebra proofs with their nonzero conditions.

Clause~3 implies clause~5 since every tuple in every model
realizes an isolated type.
Conversely, realize any specified \(p\) in a countable
model by \cref{thm:p03-rn-type-realization}.
Under clause~5 the realizing tuple is atomic, isolating \(p\).
Clause~3 implies clause~6 by the isolated-extension lemma,
and clause~6 implies clause~7.

Suppose clause~7 holds. Take countably infinite models
\(\mathcal A,\mathcal B\) and enumerate their domains.
Start with the empty elementary map: both models agree
on all sentences because \(T\) is complete.
For a forth step, take the first source element \(a\)
outside the finite current domain.
Its complete type over that domain translates via the
elementary map into a complete type over the finite range.
This translated type is consistent with the full named
theory of \(\mathcal B\): a finite conjunction of its
formulas has an existential closure true over the domain,
hence, by elementarity, true over the range.
Compactness then gives consistency of the entire named type.
Finite-parameter saturation supplies a realization \(b\)
in \(\mathcal B\).
The type contains the inequalities to each old domain
element, so \(b\) is outside the old range.
Adding \(a\mapsto b\) preserves every formula at the
enlarged tuple and thus extends the elementary map.
At a back step apply the same argument to the inverse
map and a new target element, using saturation of
\(\mathcal A\).
Alternate these two steps. The enumeration rule includes
every element in the eventual domain and range.
The union is an onto elementary bijection: each formula's
finitely many arguments already occur at a finite stage.
Atomic formulas ensure preservation of all relations,
constants and operations, so this is an isomorphism.
This proves clause~7 implies clause~1.

Under the first seven conditions, choose any countably
infinite model, whose existence follows from the
countable construction. It is atomic, hence prime by
\cref{thm:p03-vaught-prime-atomic}, and is saturated
by clause~6. This proves clause~8.
Conversely a prime model in clause~8 is atomic by that
same proved equivalence. Its saturation over the empty
set realizes every complete finite-tuple type; atomicity
isolates each realized type. Clause~3 follows.

For countable languages with functions, use the
domain-preserving graph presentation of
\cref{thm:p03-vaught-functional-bridge}.
Its translated theory is complete, satisfiable and
has only infinite models.
The two-way formula translation preserves equivalence
classes, complete types and isolators.
The mutually inverse model correspondence preserves
isomorphisms, elementary embeddings, atomicity,
primeness and finite-parameter saturation.
Naming parameters is unchanged.
All clauses and implications therefore transfer back
without any finite-signature restriction.
\end{proof}
\end{annexitem}

\begin{annexitem}{annex:p3.rn.orbits}{thm:p03-rn-orbits}{Back-and-forth starting from the prescribed tuple map}
\begin{proof}
Automorphisms preserve complete types.
Conversely equal-type tuples specify a well-defined
bijection between their finite entry sets: equality
formulas identify exactly the same repeated coordinates.
The common complete type makes this map elementary.
Start the preceding finite-parameter back-and-forth in
\(\mathcal M\) and itself with this particular map.
Every stage extends it, so the final onto elementary
map is an automorphism with the prescribed tuple image.
Saturation over the empty set realizes every complete
type in \(\mathcal M\).
The two implications thus identify types with orbits.
For zero-tuples both sets have their one empty member.
\end{proof}
\end{annexitem}

\begin{annexitem}{annex:p3.rn.finite-boundary}{thm:p03-rn-finite-boundary}{Finite-model uniqueness and both cardinal bounds}
\begin{proof}
The supplied finite complete-theory proof
(\cref{thm:p03-vaught-finite-boundary}) uses a finite
intersection argument on the possible domain bijections
to show that all models are isomorphic for the whole
countable relational/constant signature, not only for
each finite reduct separately.
For a function signature apply the graph presentation
of \cref{thm:p03-vaught-functional-bridge}; the model
correspondence preserves its finite cardinality and
isomorphism in both directions.
This proves the same uniqueness in the allowed setting.

Compactness with coordinates named realizes every
consistent complete \(n\)-type in a \(T\)-model.
All such models are isomorphic to the finite model
\(\mathcal A\), so the type is realized among its
\(m^n\) tuples. This proves the first bound.
Evaluation on \(\mathcal A^n\) injects formula equivalence
classes into its power set: equal truth sets imply
agreement in every isomorphic \(T\)-model, hence
equivalence modulo \(T\). This proves the second bound.
For \(n=0\), the set of tuples is the singleton empty
tuple, giving bounds \(1\) and \(2\).
The proved finite case supplies isolation, atomicity,
primeness and saturation; these transfer through the
graph presentation as well.
The exact-cardinality sentence in \(T\) excludes
infinite models. At-most-countable categoricity holds,
but the existence clause for \(\omega\)-categoricity fails.
\end{proof}
\end{annexitem}

\begin{annexitem}{annex:p3.rn.incomplete-boundary}{thm:p03-rn-incomplete-boundary}{Exact-cardinality types in the empty equality theory}
\begin{proof}
In the equality-only language each domain bijection is
an isomorphism. All countably infinite models are
therefore isomorphic, and a set of every positive
finite size is also a model of the empty theory.
For \(m\geq1\), let \(\rho_m\) say that the universe
has exactly \(m\) elements. Choose an element in a
model of that size; its complete one-type contains
\(\rho_m\). In a model of size \(j\ne m\), every
one-type contains \(\neg\rho_m\).
These types are distinct for all positive \(m\).
Thus \(S_1(T)\) is infinite.
The empty theory decides none of these mutually
exclusive cardinality alternatives: completeness
and the infinite-only standing have been removed.
\end{proof}
\end{annexitem}

\begin{annexitem}{annex:p3.rn.uncountable-language}{thm:p03-rn-uncountable-language}{A continuum of images at every unnamed point}
\begin{proof}
Take the specified structure and its complete theory.
The formulas \(x=c_i\) isolate distinct consistent
one-types, since \(T\) contains all
\(c_i\ne c_j\) for \(i\ne j\).
So \(S_1(T)\) is infinite.
There are \(2^{\aleph_0}\) functions \(\omega\to\omega\),
making the signature uncountable.

Choose an injection \(e:2^{<\omega}\to\omega\)
from finite binary strings.
For each infinite binary sequence \(\eta\in2^\omega\)
define \(g_\eta(i)=e(\eta\upharpoonright i)\).
For \(\eta\ne\nu\), their initial strings agree
at only finitely many lengths. Thus the set
\[
 E_{\eta,\nu}=\{i:g_\eta(i)=g_\nu(i)\}
\]
is finite. The original structure satisfies, and its
theory therefore contains,
\[
 \forall x\left(f_{g_\eta}(x)=f_{g_\nu}(x)\to
                  \bigvee_{i\in E_{\eta,\nu}}x=c_i\right).
\]
An empty disjunction means \(\bot\).
In any \(\mathcal N\models T\), a point \(a\) outside
all named constants consequently has pairwise distinct
images \(f_{g_\eta}^{\mathcal N}(a)\).
The map sending \(\eta\) to that image is an injection
of \(2^\omega\) into the domain, so the domain has at
least continuum many elements.

Every at most countable model thus consists exactly
of its pairwise distinct \(c_i\)'s.
The sentences \(f_g(c_i)=c_{g(i)}\) in \(T\), for every
function \(g\) and index \(i\), make
\(i\mapsto c_i^{\mathcal N}\) an isomorphism from the
original structure onto such a model.
There is no finite model because all constants are distinct.
This proves the required existence and uniqueness.

There also is a consistent complete type outside all
named points. Each finite subset of
\(\{x\ne c_i:i<\omega\}\) is realized in the original
structure. Compactness and extension of a consistent
type to a complete one provide such a type.
Every individual formula still uses only finitely many
symbols; first-order compactness is valid for this
uncountable language. It is the countable-language
omission implication, not compactness, that cannot
be applied here.
\end{proof}
\end{annexitem}

\begin{annexitem}{annex:p3.rn.dlo-extension}{thm:p03-rn-dlo-extension}{The equality, interval and two end-ray cases}
\begin{proof}
Write the distinct old source points as
\(a_1<\cdots<a_k\), and their images as
\(b_1<\cdots<b_k\).
An old source point keeps its existing image.
If the source point is new and \(k=0\), choose any
target point from the nonempty target domain.
For \(k>0\), it lies below \(a_1\), above \(a_k\),
or between a unique adjacent pair.
Use respectively the no-left-endpoint axiom,
the no-right-endpoint axiom, or density to choose
a target in the corresponding cut.
The chosen point is outside the old range and has
exactly the required comparisons, so the enlarged
map is an order-preserving finite bijection.
Apply the same construction to the inverse map for
the back direction.

Repeated density between an interval endpoint and
a selected interior point gives any prescribed
finite number of distinct interior points.
Repeated no-endpoint selection gives arbitrarily
many points in either end ray.
No finite forbidden set exhausts any of these cuts,
which justifies a fresh target whenever required.
\end{proof}
\end{annexitem}

\begin{annexitem}{annex:p3.rn.dlo-pattern-truth}{thm:p03-rn-dlo-pattern-truth}{Formula induction before quantifier elimination}
\begin{proof}
Equal order patterns prescribe a well-defined
order-preserving finite bijection between tuple entry
sets, including all repeated-coordinate identifications.
Prove by induction on formulas that these correspondences
preserve truth between any two models.
Equality and order atoms are preserved by definition.
Boolean operations follow from the induction hypothesis.
For \(\exists y\,\varphi(\bar x,y)\), choose a source
witness and extend the finite map to it by the proved
extension lemma. The extended tuples have the same
pattern; the induction hypothesis transfers \(\varphi\),
and hence the existential formula.
The inverse extension proves the reverse implication.
Universal quantification follows by negation.
Bound-variable renaming avoids capture and covers formulas
with unused free slots.
This proves all formulas without any countability or
quantifier-elimination assumption.

The ordered rationals are a model, so the theory is
satisfiable. The empty-tuple case shows that every two
models agree on all sentences; the truth value in the
rational model decides each sentence in all models.
Thus the theory is complete.

To realize an ordered partition with \(k\) blocks,
choose \(k\) increasing distinct points, using
nonemptiness and successive no-right-endpoint choices;
assign each point to the coordinates of its block.
For \(n=0\) use the empty tuple.
There are finitely many patterns at each fixed \(n\).
For any formula retain precisely the patterns on
which it is true. The proved induction makes this
truth independent of model and chosen realizing tuple.
Every tuple has one and only one pattern, so the
formula is equivalent in every model to the disjunction
of the retained patterns.
Use \(\bot\) when none are retained; for zero-tuples
use \(\top\) or \(\bot\).
These are quantifier-free formulas, establishing
quantifier elimination as a consequence.
\end{proof}
\end{annexitem}

\begin{annexitem}{annex:p3.rn.dlo-type-count}{thm:p03-rn-dlo-type-count}{The ordered-block count and the full Boolean quotient}
\begin{proof}
Each pattern is consistent and decides every formula
by the preceding theorem, so it isolates a complete
type. Distinct patterns disagree on an equality or
order atom and yield distinct types.
Every complete type has a realization by the
countable-realization lemma. Its tuple has one
of these patterns, so there are no additional types.

Choose an unordered partition of the \(n\) indices
into \(k\) blocks in \({n\brace k}\) ways, then order
its distinct blocks in \(k!\) ways.
Sum over \(k\) for the stated count.
For \(n=0\), only \(k=0\) contributes, with value \(1\).
For \(n=1\), only one block contributes, with value \(1\).
For \(n=2\), the partition counts are \(1\) and \(1\), giving
\(1+2=3\).
For \(n=3\), they are \(1,3,1\) for one, two and three
blocks, giving \(1+6+6=13\).

Formula equivalence classes correspond to subsets
of the finite pattern set. Every subset occurs as
the disjunction of its isolators, and distinct
subsets differ on a realized pattern, hence give
inequivalent formulas.
This is a bijection with the full power set,
giving the exact Boolean-quotient size.
\end{proof}
\end{annexitem}

\begin{annexitem}{annex:p3.rn.dlo-uniqueness}{thm:p03-rn-dlo-uniqueness}{An onto order map that extends its initial assignment}
\begin{proof}
The rationals are countably infinite. Their strict
order is linear, \((a+b)/2\) lies between \(a<b\),
and \(a-1,a+1\) lie on either side of each \(a\).
Thus they satisfy the axioms.
The endpoint and density axioms likewise make
every model infinite.

Enumerate both countable model domains and start
with the empty map or the prescribed finite order map.
At each even stage extend to the first source point
outside the current domain, using the extension lemma.
At each odd stage use the inverse map and the first
target point outside the range.
The maps are nested finite order bijections.
The enumeration rule eventually includes every
point in the domain and range.
Their union is a bijection, and each pair of points
occurs at a finite stage where order is preserved
in both directions.
It is therefore an order isomorphism and extends
the initial map.
Taking one model to be the rationals proves
uniqueness; taking both to be the rationals proves
homogeneity.
The proved pattern/type identification yields the
orbit statement, and existence plus uniqueness yields
\(\omega\)-categoricity directly.
Both endpoint clauses were used in the two end-ray
steps of the finite extension.
\end{proof}
\end{annexitem}

\begin{annexitem}{annex:p3.rn.relational-application}{thm:p03-rn-relational-application}{Finite atomic patterns and genuine map extension}
\begin{proof}
At each fixed \(n\), there are finitely many equality
atoms and finitely many relation atoms on the \(n\)
slots: the signature and each relation's arity
are finite. There are finitely many conjunctions
assigning truth or falsity to every such atom.
Retain those realized in \(\mathcal M\).
Two tuples with the same retained conjunction prescribe
a well-defined bijection of their finite entry sets,
by the equality atoms.
The relation atoms make it an isomorphism of the
finite induced substructures, including relation
instances with repeated entries.
Homogeneity extends it to an automorphism.
Consequently the tuples agree on every formula.

For any formula, choose the finitely many realized
patterns making it true.
The universal closure of its equivalence to their
disjunction is true in \(\mathcal M\), hence belongs
to \(\operatorname{Th}(\mathcal M)\).
This proves quantifier elimination.
The disjunction of all retained patterns is likewise
exhaustive in every model of that complete theory.
Each retained pattern is consistent and, by the
formula equivalences, decides every formula, so
it isolates a complete type.
A complete consistent type contains one retained
pattern: it cannot contain the negations of all
disjuncts in a finite exhaustive disjunction.
Thus there are finitely many complete types in
each arity; the zero-type is unique by completeness.
The finite-type criterion proves categoricity.
Absence of functions is used in the finite atom
list and in treating the finite entry set as
a closed substructure.
\end{proof}
\end{annexitem}

\begin{annexitem}{annex:p3.rn.function-boundary}{thm:p03-rn-function-boundary}{The unbounded term list in a finite signature}
\begin{proof}
A nonempty subset of \(\mathbb Z\) closed under
both \(s,t\) contains every integer successor
and predecessor of each of its elements.
It is all of \(\mathbb Z\).
Thus every nonempty finitely generated substructure
is the whole structure; every isomorphism between
such substructures is already an automorphism,
and there is no nonempty finite closed substructure.
For each \(k\geq0\), the pair \((0,k)\) satisfies
\(y=s^k(x)\), whereas \((0,j)\) for \(j\ne k\)
does not. These complete two-types are all distinct.
The complete theory has a countable language and
an infinite model, so its infinite two-type space
violates the necessary criterion for
\(\omega\)-categoricity.
The terms \(s^k(x)\) give infinitely many atomic
formulas despite the two-symbol signature.
That is exactly the finite-pattern step that
the relational application would lose.
\end{proof}
\end{annexitem}
